\documentclass{report}
\usepackage{amsmath,amssymb}
\setlength{\parindent}{0mm}
\setlength{\parskip}{1em}
\begin{document}
\begin{center}
\rule{6in}{1pt} \
{\large Hillary Fairbanks \\
{\bf High Dimensional Uncertainty Quantification via Multilevel Monte Carlo}}

1859 23rd st \\ APT C \\ Boulder \\ CO \\ 80302
\\
{\tt hillary.fairbanks@colorado.edu}\\
Alireza Doostan\\
Christian Ketelsen\end{center}

Multilevel Monte Carlo (MLMC) has been shown to be a cost effective way
to compute moments of desired quantities of interest in stochastic
partial differential equations when the uncertainty in the data is
high-dimensional. In this talk, we investigate the improved performance
of MLMC versus standard Monte Carlo. While both methods converge at a
rate independent of the dimension of the stochastic input, the use of a
series of nested grids in MLMC allows us to improve the convergence by
allocating more computational work onto the coarse grids, while the
costly fine grid solves require fewer samples. In addition, we find that
further variance reduction techniques have promise in speeding up
convergence to a solution. As an application, we consider a thermally
driven flow problem with random fluctuations in temperature along one
wall of a square domain. Given the random input of the initial
temperature, we wish to determine the expected value of the final mean
heat flux along the opposite wall.


\end{document}
